\documentclass[11pt,a4paper]{article}

% ---------- fontes e língua ----------
\usepackage[T1]{fontenc}
\usepackage[utf8]{inputenc}
\usepackage{lmodern}
\usepackage{tgpagella}
\usepackage{tgheros}
\usepackage[portuguese]{babel}
\usepackage{microtype}

% ---------- página e cores ----------
\usepackage[a4paper,top=2.6cm,bottom=2.4cm,left=2.4cm,right=2.4cm]{geometry}
\usepackage[table]{xcolor}
\definecolor{azul}{RGB}{28,62,112}
\definecolor{azulclaro}{RGB}{226,234,246}
\definecolor{cinza}{RGB}{95,102,115}

% ---------- tabelas, listas, caixas ----------
\usepackage{booktabs}
\usepackage{array}
\usepackage{tabularx}
\usepackage{enumitem}
\usepackage{titlesec}
\usepackage{fancyhdr}
\usepackage{lastpage}
\usepackage{tcolorbox}

\setlength{\parindent}{0pt}
\setlength{\parskip}{6pt}

\setlist{nosep, leftmargin=*, itemsep=3pt}
\setlist[itemize]{label=\textcolor{azul}{\textbullet}}
\setlist[enumerate]{label=\textcolor{azul}{\sffamily\bfseries\arabic*.}}

\titleformat{\section}{\sffamily\Large\bfseries\color{azul}}{\thesection}{0.7em}{}[{\color{azul}\titlerule[0.6pt]}]
\titleformat{\subsection}{\sffamily\large\bfseries\color{azul}}{\thesubsection}{0.6em}{}
\titlespacing*{\section}{0pt}{1.6em}{0.8em}
\titlespacing*{\subsection}{0pt}{1.1em}{0.3em}

\pagestyle{fancy}
\setlength{\headheight}{15pt}
\fancyhf{}
\renewcommand{\headrule}{{\color{azul}\hrule width\headwidth height\headrulewidth}}
\renewcommand{\headrulewidth}{0.4pt}
\fancyhead[L]{\sffamily\small\color{cinza} Universidade da Beira Interior}
\fancyhead[R]{\sffamily\small\color{cinza} Engenharia de Software \textperiodcentered{} 2026/27}
\fancyfoot[L]{\sffamily\small\color{cinza} Análise do Produto \textperiodcentered{} Sprint 1}
\fancyfoot[R]{\sffamily\small\color{cinza} Página \thepage\ de \pageref{LastPage}}

\fancypagestyle{primeira}{%
  \fancyhf{}%
  \renewcommand{\headrulewidth}{0pt}%
  \fancyfoot[L]{\sffamily\small\color{cinza} Análise do Produto \textperiodcentered{} Sprint 1}%
  \fancyfoot[R]{\sffamily\small\color{cinza} Página \thepage\ de \pageref{LastPage}}%
}

\newcolumntype{Y}{>{\raggedright\arraybackslash}X}
\newcolumntype{L}[1]{>{\raggedright\arraybackslash}p{#1}}
\newcolumntype{N}[1]{>{\raggedright\arraybackslash\sffamily\bfseries}p{#1}}
\newcommand{\th}[1]{\textcolor{white}{\sffamily\bfseries #1}}

\tcbset{
  caixa/.style={colback=white, colframe=azul, coltitle=white, colbacktitle=azul,
    fonttitle=\sffamily\bfseries, arc=3pt, boxrule=0.6pt,
    left=9pt, right=9pt, top=6pt, bottom=6pt}
}

\begin{document}
\thispagestyle{primeira}

% ---------- cabeçalho do documento ----------
\begin{tcolorbox}[colback=azul, colframe=azul, coltext=white, arc=3pt, boxrule=0pt,
    left=16pt, right=16pt, top=14pt, bottom=14pt]
  \sffamily
  {\footnotesize\bfseries\MakeUppercase{Universidade da Beira Interior}}\\[2pt]
  {\footnotesize Faculdade de Engenharia \textperiodcentered{} Departamento de Informática}\\[16pt]
  {\fontsize{28}{32}\selectfont\bfseries Análise do Produto}\\[4pt]
  {\large Aplicação Universitária UBI \textperiodcentered{} Sprint 1 \textperiodcentered{} 28/09}
\end{tcolorbox}

%--------------------------------------------------------------------
\section{Visão e âmbito}

A informação e os serviços da universidade estão espalhados por vários sítios, o que complica o dia a dia e a orientação no campus. A ideia é juntar tudo numa só aplicação, pensada para estudantes, docentes e funcionários.

Para já, o âmbito inicial (ainda por confirmar) é este:

\begin{itemize}
  \item consultar horários e salas;
  \item receber avisos quando há alterações;
  \item permitir a comunicação entre docentes e estudantes.
\end{itemize}

As restantes funcionalidades ficam para versões futuras, conforme as prioridades, as dependências e a opinião de quem usar a aplicação.

%--------------------------------------------------------------------
\section{Stakeholders}

\begin{itemize}
  \item \textbf{Estudantes:} são os utilizadores principais. Consultam horários, notas, refeitório e transportes.
  \item \textbf{Docentes:} dão as aulas, tratam das avaliações e falam com as turmas.
  \item \textbf{Funcionários administrativos:} tratam das inscrições, das propinas e dos processos administrativos.
  \item \textbf{Serviços Académicos (SASUBI/Reitoria):} responsáveis pelos processos e pelas regras de negócio.
  \item \textbf{Equipa de TI da UBI:} manutenção, suporte técnico e integração.
  \item \textbf{Serviços de Ação Social:} bolsas, refeitório e residências.
  \item \textbf{Segurança/Portaria:} controlo de acessos aos edifícios, se for preciso.
  \item \textbf{Fornecedores externos:} empresas que podem ser integradas, como a cantina e os transportes.
  \item \textbf{Visitantes/Candidatos:} acesso limitado à informação pública e às candidaturas.
\end{itemize}

%--------------------------------------------------------------------
\section{Funcionalidades identificadas}

\begin{enumerate}
  \item Consulta de horários e salas
  \item Consulta de notas
  \item Consulta de propinas
  \item Notificações sobre alterações de aulas e frequências
  \item Reservas de espaços (salas de estudo, anfiteatros, entre outros)
  \item Mapa interativo da universidade, com salas e serviços da UBI
  \item Comunicação entre docentes e estudantes
  \item Gestão de pedidos de documentos
\end{enumerate}

%--------------------------------------------------------------------
\section{Personas}

\begin{tcolorbox}[caixa, title={Inês Carvalho \textperiodcentered{} estudante}]
  19 anos, frequenta o 1.\textordmasculine{} ano de Engenharia Informática. É de Viseu, estuda na Covilhã e vive numa residência. Usa o telemóvel para quase tudo.

  Quer ver o horário, as salas e as notas num só sítio, consultar o menu da cantina e orientar-se pelos vários polos. Frustra-se por andar de site em site e por não saber a quem recorrer nos assuntos administrativos.
\end{tcolorbox}

\begin{tcolorbox}[caixa, title={Professora Helena Duarte \textperiodcentered{} docente}]
  52 anos, dá aulas de Matemática há mais de 20. Tem conhecimentos tecnológicos medianos e prefere ferramentas simples.

  Quer comunicar depressa com os alunos, lançar sumários e notas, e reservar salas e equipamentos. Frustra-se com tarefas repetitivas e com tantos sites diferentes.
\end{tcolorbox}

\begin{tcolorbox}[caixa, title={Paula Rodrigues \textperiodcentered{} funcionária dos Serviços Académicos}]
  45 anos, atende estudantes presencialmente e por email. O volume de pedidos sobe no início do semestre e nas inscrições para exames.

  Quer reduzir os pedidos repetitivos e tratar dos processos de forma digital. Frustra-se com as filas, com os pedidos incompletos e com perguntas cuja resposta já está publicada.
\end{tcolorbox}

\begin{tcolorbox}[caixa, title={João Almeida \textperiodcentered{} estudante com mobilidade reduzida}]
  20 anos, estuda Marketing e desloca-se em cadeira de rodas.

  Quer consultar percursos e edifícios acessíveis, reservar apoio e receber avisos quando um elevador está avariado. Frustra-se quando só descobre no local que um acesso não está disponível.
\end{tcolorbox}

%--------------------------------------------------------------------
\section{Cenários de utilização}

\subsection{Cenário 1: consultar o horário e chegar à aula}

Todas as manhãs a Inês acorda sem saber se o horário mudou. Abre a aplicação e vê as aulas do dia, com horas e salas, e confirma se há alguma notificação de alteração antes de sair de casa. Já na universidade, vê onde fica a sala e segue para a aula. Se a informação estiver desatualizada ou a sala não estiver identificada, a aplicação avisa que não tem os dados e indica um contacto.

\subsection{Cenário 2: avisar a turma de uma alteração}

A professora Helena tem de avisar a turma de que a aula mudou de sala. Na funcionalidade de comunicação escolhe a turma e a aula, escreve a alteração e confirma o envio. Os estudantes recebem uma notificação com o essencial e a Helena consegue ver que a mensagem foi enviada. Se não tiver autorização, ou se a turma não estiver disponível, o sistema explica o problema e não envia nada.

\subsection{Cenário 3: pedir um documento académico}

A Paula recebe muitos pedidos de declarações e alguns chegam sem os dados necessários. Na aplicação, o estudante vê os documentos disponíveis, escolhe o que quer e preenche os campos obrigatórios. Antes de enviar, revê o pedido. O sistema confirma a submissão e mostra um número ou um estado para acompanhar o pedido. Se faltar algum dado, assinala os campos a corrigir.

\subsection{Cenário 4: encontrar um percurso acessível}

O João tem de chegar a um edifício num polo que não conhece. Abre o mapa, procura o edifício e escolhe a opção de percurso acessível. A aplicação mostra um percurso com entradas acessíveis. Se um elevador estiver avariado, o aviso aparece bem à vista e, quando existir, é apresentada uma alternativa. Se não existir, a aplicação diz que não consegue garantir um percurso acessível.

\subsection{Cenário 5: consultar uma nota ou as propinas}

A Inês quer confirmar uma nota que saiu ou ver informação sobre as propinas. Abre a secção certa e vê os dados da sua conta. A aplicação separa o que é atual do que não está disponível ou ainda não foi publicado, e nunca mostra dados de outros estudantes.

%--------------------------------------------------------------------
\section{User stories iniciais}

\small
\renewcommand{\arraystretch}{1.5}
\begin{tabularx}{\textwidth}{N{1.3cm} L{3.0cm} Y}
  \rowcolor{azul} \th{ID} & \th{Tema} & \th{História} \\
  \rowcolor{azulclaro!55}
  US01 & Horário &
    Como estudante, quero ver o meu horário e as salas, para planear o dia sem andar de site em site. \\
  \rowcolor{white}
  US02 & Notificações &
    Como estudante, quero receber avisos quando uma aula ou uma frequência muda, para agir com informação atualizada. \\
  \rowcolor{azulclaro!55}
  US03 & Comunicação &
    Como docente, quero avisar a turma de uma alteração, para os estudantes ficarem a saber a tempo. \\
  \rowcolor{white}
  US04 & Percurso acessível &
    Como estudante com mobilidade reduzida, quero ver percursos acessíveis e avisos sobre elevadores, para planear a deslocação. \\
  \rowcolor{azulclaro!55}
  US05 & Pedido de documento &
    Como estudante, quero submeter um pedido de documento e acompanhá-lo, para evitar deslocações e saber em que estado está. \\
  \rowcolor{white}
  US06 & Gestão de pedidos &
    Como funcionária dos Serviços Académicos, quero receber os pedidos com os campos necessários, para haver menos pedidos incompletos. \\
  \bottomrule
\end{tabularx}
\renewcommand{\arraystretch}{1}
\normalsize

São histórias candidatas ao Product Backlog. Ainda falta refiná-las, dar-lhes prioridade e escolher as que cabem no sprint, com critérios de aceitação combinados.

%--------------------------------------------------------------------
\section{Priorização preliminar}

\small
\renewcommand{\arraystretch}{1.5}
\begin{tabularx}{\textwidth}{N{0.9cm} L{4.6cm} L{1.9cm} Y}
  \rowcolor{azul} \th{N.\textordmasculine} & \th{Funcionalidade} & \th{Prioridade} & \th{Justificação} \\
  \rowcolor{azulclaro!55}
  1 & Horários e salas & Alta & Usa-se todos os dias e é informação essencial para os estudantes. \\
  \rowcolor{white}
  2 & Notas & Alta & Ajuda a acompanhar o semestre. \\
  \rowcolor{azulclaro!55}
  3 & Propinas & Média & É importante, mas consulta-se menos vezes. \\
  \rowcolor{white}
  4 & Notificações de alterações & Alta & Ajuda a reagir às mudanças e evita deslocações ou faltas. \\
  \rowcolor{azulclaro!55}
  5 & Reservas de espaços & Baixa & É útil, mas a prioridade tem de ser validada com os utilizadores e com o âmbito. \\
  \rowcolor{white}
  6 & Mapa interativo & Baixa & Pode ser mais valioso para os novos alunos, os visitantes e as pessoas com necessidades de acessibilidade. \\
  \rowcolor{azulclaro!55}
  7 & Comunicação docente--estudante & Alta & Facilita a comunicação direta sobre aulas e atividades. \\
  \rowcolor{white}
  8 & Pedidos de documentos & Média & Pode reduzir o trabalho administrativo e os pedidos repetitivos. \\
  \bottomrule
\end{tabularx}
\renewcommand{\arraystretch}{1}
\normalsize

%--------------------------------------------------------------------
\section{Por confirmar}

\begin{itemize}
  \item O âmbito, o nome oficial do produto e os utilizadores-alvo. E também se a aplicação se liga a sistemas que já existem ou se só junta ligações e informação.
  \item As personas e os cenários têm de ser validados com estudantes, docentes e funcionários. Até lá são suposições nossas (proto-personas), não factos.
  \item Como vamos obter os horários, as notas, as propinas, os dados de acessibilidade e o estado dos elevadores, e quem mantém esses dados.
  \item As prioridades: o que é muito importante para poucas pessoas (como a acessibilidade) não deve ficar para trás só porque a maioria o usa menos.
  \item Os critérios de aceitação das histórias escolhidas, ligados às funcionalidades. Não vamos pôr tudo no Sprint 1.
\end{itemize}

\end{document}
